\section{The four-branch foundation}\label{sec:four-foundation}
The five-branch argument uses the following four-branch result whenever
a reduction removes one branch vertex.
\begin{theorem}\label{thm:four}
For every nonempty finite tree $U$ with $b(U)\le4$, there is an
orientation $D=D(U)$ of all its edges such that \eqref{eq:tas} holds for
every nonempty finite loopless tournament.
\end{theorem}
\begin{theorem}\label{thm:four-kernel}
The orientation in Theorem~\ref{thm:four} can be chosen so that
$Z_D(K,\mu)\le1$ for every finite probability space and every real skew
kernel bounded in absolute value by one.
\end{theorem}
Appendix~\ref{sec:three} proves the assertion for at most three branch
vertices. Section~\ref{sec:geometry} reduces to the two skeletons,
Section~\ref{sec:terminals} supplies the terminal estimates and the
internal side tables; Sections~\ref{sec:chain} and~\ref{sec:star}
then close the chain and the star. These inputs also serve the five-branch
case.
